\lecture{3.5}{Assumptions for Identifying the ATE}
\label{lec:3-5}

\subsection{Consistency}

To identify the ATE, the observed outcome must be linked to the relevant potential outcome.

\begin{assumptionbox}
\textbf{Consistency.} If a unit is observed with $A=a$, then
\[
Y=\Ya{a}.
\]
For binary $A$,
\[
Y=A\Yone+(1-A)\Yzero.
\]
\end{assumptionbox}

Consistency requires that the treatment versions represented by $A=a$ are well-defined enough that the observed treatment corresponds to the intervention used to define $\Ya{a}$.

\subsection{Conditional Randomization}

The second assumption is conditional randomization, also called conditional exchangeability or conditional ignorability.

\begin{assumptionbox}
\textbf{Conditional randomization.} For $a\in\{0,1\}$,
\[
\Ya{a}\ci A\given W.
\]
Within strata of measured covariates $W$, the observed treatment carries no information about the potential outcome under treatment level $a$.
\end{assumptionbox}

This assumption is usually the hardest to justify in observational studies. It says that $W$ contains enough information to control confounding.

\subsection{Positivity}

Even if treatment is as-if randomized within covariate strata, we must have observed treatment variation within those strata.

\begin{assumptionbox}
\textbf{Positivity.} For each treatment level $a\in\{0,1\}$,
\[
\Prob(A=a\given W=w)>0
\]
for every covariate value $w$ with $\Prob(W=w)>0$.
\end{assumptionbox}

In words, every covariate subgroup in the target population must have some chance of receiving each treatment level. If no mildly ill patients ever receive a risky experimental treatment, then the observed data contain no direct information about outcomes under that treatment for mildly ill patients.

\subsection{Why Positivity Is Partly Empirically Checkable}

Unlike exchangeability, positivity is stated in terms of the observed data distribution. It involves $\Prob(A=a\given W=w)$, not counterfactual outcomes. Therefore it can be partly checked.

For discrete $W$, one can inspect whether each covariate stratum contains both treated and untreated observations. For continuous or high-dimensional $W$, exact strata may be sparse, so analysts examine estimated propensity scores
\[
g_a(W)=\Prob(A=a\given W)
\]
and check whether they are close to zero.

Empirical checks are not perfect. A finite sample may fail to show a treatment level in a subgroup even if the population probability is positive. Conversely, a few observations in a subgroup do not guarantee enough information for stable estimation. Still, positivity has observable implications.

\subsection{Why Positivity Matters for Observed Conditional Means}

The G-computation formula will involve terms such as
\[
\E(Y\given A=a,W=w).
\]
This conditional mean is defined from the observed data distribution only when $\Prob(A=a,W=w)>0$. Since
\[
\Prob(A=a,W=w)=\Prob(A=a\given W=w)\Prob(W=w),
\]
if $\Prob(W=w)>0$ but $\Prob(A=a\given W=w)=0$, then the subgroup $(A=a,W=w)$ does not exist in the observed population. The conditional mean cannot be learned from observed outcomes.

\begin{keyidea}
Exchangeability says treated and untreated people are comparable within $W$ strata. Positivity says both treated and untreated people actually exist within those strata.
\end{keyidea}
